Solución FEM acoplada calculada para P5 con elementos triangulares
cuadráticos \(P_2\) y la malla fina alineada con la interfaz. (a) Campo
vertical \(d_y\) en micrómetros; (b) potencial \(\phi\) en voltios; (c)
contorno exterior e interfaz originales y deformados, cuya escala visual es el factor calculado
28.3071; (d) cambio de la flecha de punta y de la carga libre del
electrodo superior al refinar. La deformada se dibuja con coordenadas
geométricas \(x,y\) en metros y con razón de aspecto ajustada para facilitar
la lectura. Como \(L/H=100\), los ejes no usan la misma escala gráfica; además,
el desplazamiento se multiplica por el factor indicado solo para ser visible.
La carga mostrada usa
\(Q_T=-\int_{\Gamma_T}D\!\cdot n\,ds\), mientras el flujo saliente del
material tiene el signo contrario.
